\chapter{Milne Theorem 3.2: Rational maps into abelian varieties}
\label{chap:Albanese_Thm32RationalMapExtension}
% archon:covers AlgebraicJacobian/Albanese/Thm32RationalMapExtension.lean

% STRATEGY NOTE (iter-174). This chapter is the A.4.c sub-row of Route A
% (the positive-genus arm of \texttt{nonempty\_jacobianWitness}). It is the
% join of the two preceding sub-rows: A.4.a, the codim-1 indeterminacy /
% codim-\(\geq 2\) extension package (\cref{chap:Albanese_CodimOneExtension}),
% and A.4.b, the Auslander--Buchsbaum corollary
% ``regular local \(\Rightarrow\) Cohen--Macaulay''
% (\cref{chap:Albanese_AuslanderBuchsbaum}). The single output of this
% chapter is Milne's Theorem~3.2 of Chapter~I~\S 3 of \emph{Abelian Varieties}:
% \emph{every rational map from a nonsingular variety to an abelian variety
% extends, uniquely, to a regular morphism}. This is the input that A.4.d
% (\cref{chap:Albanese_AlbaneseUP}) consumes to upgrade the symmetric-product
% rational map \(C^{(g)} \dashrightarrow A\) to a regular morphism \(J \to A\),
% which in turn yields the Albanese universal property of \(\Pic^0_{C/k}\)
% (Milne III~\S 6, Proposition~6.1).
%
% This chapter is also the canonical home of the Lean target
% \texttt{AlgebraicGeometry.Scheme.RationalMap.extend\_to\_av}. A placeholder
% lemma block \texttt{lem:rational\_map\_to\_av\_extends} was reserved
% earlier in \cref{chap:AbelianVarietyRigidity} (under the alternative Lean
% name \texttt{AlgebraicGeometry.rationalMap\_to\_av\_extends}); a Plan-Agent
% reconciliation pass should retarget that reservation at the theorem proved
% here (see "Notes for Plan Agent" in the writer report).
%
% The Lean target file \texttt{AlgebraicJacobian/Albanese/Thm32RationalMapExtension.lean}
% does not yet exist; this chapter is the prover-ready specification for
% creating it.

\section{Setup and statement}
\label{sec:thm32_setup}

Fix an algebraically closed field \(\bar k\), a smooth variety \(X\) over \(\bar k\) (in the
project's convention: a geometrically integral, geometrically reduced, separated scheme of
finite type over \(\bar k\) that is moreover smooth, hence nonsingular at every closed point),
and an abelian variety \(A\) over \(\bar k\) (in particular \(A\) is complete, smooth, and a
group object in \(\mathrm{Sch}/\bar k\)). A rational map
\(f \colon X \dashrightarrow A\) is, in the project's notation, the equivalence class of a
regular morphism \(\varphi_U \colon U \to A\) defined on a dense open subset
\(U \subseteq X\), two such pairs \((U, \varphi_U)\) and \((U', \varphi_{U'})\) being equivalent
if \(\varphi_U\) and \(\varphi_{U'}\) agree on \(U \cap U'\).

The result of this chapter is:

\begin{theorem}

\leanok
[Milne Theorem~3.2: a rational map from a nonsingular variety to an abelian variety extends]
  \label{thm:rational_map_to_av_extends}
  \lean{AlgebraicGeometry.Scheme.RationalMap.extend_to_av}
  \uses{thm:codim_one_extension, lem:milne_codim1_indeterminacy, lem:av_indeterminacy_empty}

  % SOURCE: [Milne, Abelian Varieties], Theorem~3.2, \S I.3, doc p.~17
  % (read from references/abelian-varieties.pdf, PDF page 23).
  % SOURCE QUOTE:
  % "THEOREM 3.2. A rational map \(\alpha\colon V--\to A\) from a nonsingular
  %  variety to an abelian variety is defined on the whole of \(V\)."
  \textit{Source: Milne, \emph{Abelian Varieties}, Theorem~3.2, \S I.3, p.~17.}

  Let \(X\) be a nonsingular variety over an algebraically closed field \(\bar k\), let \(A\) be
  an abelian variety over \(\bar k\), and let \(f \colon X \dashrightarrow A\) be a rational
  map of \(\bar k\)-varieties (defined over \(\bar k\)).
  Then \(f\) is defined on the whole of \(X\): there is a unique regular morphism
  \(\widetilde f \colon X \to A\) whose restriction to some (equivalently, every) dense open
  set of definition \(U \subseteq X\) of \(f\) agrees with \(\varphi_U\).
\end{theorem}

The uniqueness clause is immediate from \(X\) being reduced and separated: two regular
morphisms \(X \to A\) that agree on a dense open subset agree everywhere. The substance is
existence of the extension. Milne deduces it in one line by combining the codim-\(\geq 2\)
extension result for rational maps to a complete variety with the codim-1 indeterminacy
structure result for rational maps to a group variety:

% SOURCE QUOTE PROOF:
% "PROOF. Combine Theorem 3.1 with the next lemma."
\begin{proof}
  
  Milne's one-line proof is: \emph{combine Theorem~3.1 with Lemma~3.3}.
  Theorem~3.1 (the codim-\(\geq 2\) extension theorem for a rational map from a normal
  variety to a complete variety; valuative criterion of properness) is
  \cref{thm:codim_one_extension} of the present project; Lemma~3.3 (the
  pure-codimension-\(1\) structure of the indeterminacy locus of a rational map into a
  group variety) is \cref{lem:milne_codim1_indeterminacy} of the same chapter.

  Concretely: let \(U \subseteq X\) be a maximal open set of definition of \(f\), and let
  \(Z := X \setminus U\) be its complement.
  \begin{enumerate}
    \item By \cref{thm:codim_one_extension} (Milne~3.1, the valuative-criterion extension
      to codim \(\geq 2\)), applied with the normal variety \(X\) and the complete variety
      \(A\), we have \(\mathrm{codim}_X(Z) \geq 2\).
    \item By \cref{lem:milne_codim1_indeterminacy} (Milne~3.3, the pure-codim-\(1\)
      indeterminacy structure for maps into the group variety \(A\)), the closed subset
      \(Z\) is either empty or of pure codimension \(1\) in \(X\).
    \item Both conditions together force \(Z = \emptyset\): a non-empty pure-codimension-\(1\)
      closed subset has codimension exactly \(1\), contradicting \(\mathrm{codim}_X(Z) \geq 2\).
  \end{enumerate}
  Therefore \(U = X\) and \(f\) admits a regular representative on all of \(X\), which is the
  required regular morphism \(\widetilde f \colon X \to A\). Uniqueness is the standard
  reduced-and-separated agreement principle (see the paragraph preceding the theorem).
\end{proof}

\begin{remark}
  \label{rmk:thm32_role_of_ab}
  \textbf{No depth input is needed.} The combination above runs entirely through
  the valuative criterion of properness (inside
  \cref{thm:indeterminacy_codimGe2}, applied at the DVR \(\mathcal O_{X, Z}\) of
  each codimension-\(1\) point \(Z\) --- works because \(X\) is normal and \(A\) is
  complete) and the formal extension theorem \cref{thm:codim_one_extension},
  which needs only that the indeterminacy locus is \emph{empty}. In particular
  the Auslander--Buchsbaum formula \cref{thm:auslander_buchsbaum} and its
  Cohen--Macaulay corollary \cref{cor:regular_cohen_macaulay} of
  \cref{chap:Albanese_AuslanderBuchsbaum} are \emph{not} consumed by this chain:
  a depth-\(\geq 2\) local-cohomology argument would be required only for
  Weil-style extension theorems that extend across codim-\(\geq 2\) points
  \emph{without} first emptying the locus, and codim-\(1\)-freeness alone cannot
  yield an extension for a general complete target
  (\cref{rmk:codim1_no_extension_from_codim1free}).
\end{remark}

\section{Why the proof is char-free}
\label{sec:thm32_char_free}

Each of the two inputs is char-free:
\begin{itemize}
  \item \emph{Theorem~3.1 / codim-\(\geq 2\) extension} (\cref{thm:codim_one_extension}):
    the proof uses only the valuative criterion of properness and depth-\(\geq 2\) extension
    of sections on regular local rings; both are characteristic-agnostic.
  \item \emph{Lemma~3.3 / pure-codim-\(1\) indeterminacy for maps into a group variety}
    (\cref{lem:milne_codim1_indeterminacy}): Milne's proof forms the difference map
    \(\Phi(x, y) := \varphi(x)\,\varphi(y)^{-1} \colon X \times X \dashrightarrow A\) and
    examines its zero/pole divisor along the diagonal --- a Weil-divisor argument that
    works for any algebraically closed base field. The argument uses only the group law
    on \(A\) (its smoothness as a group object), the regularity of \(X\) (so divisors and
    DVRs at codim-\(1\) points are available), and the existence of a divisor of zeros and
    poles for a nonzero rational function on a nonsingular variety. None of these inputs
    has a characteristic restriction.
\end{itemize}
Consequently \cref{thm:rational_map_to_av_extends} holds over every algebraically closed
field, in arbitrary characteristic. (Milne states this proof verbatim in characteristic-free
generality at \S I.3, pp.~16--17.)

\section{Application: the symmetric-product Albanese map (gating A.4.d)}
\label{sec:thm32_application_to_a4d}

The downstream consumer is A.4.d (\cref{chap:Albanese_AlbaneseUP}):
the Albanese universal property of the Jacobian \(J = \Pic^0_{C/k}\). Milne's proof of
Proposition~6.1 of \S III.6 (\emph{Abelian Varieties}, p.~104) constructs the unique
morphism \(\psi \colon J \to A\) factoring a given pointed regular morphism
\(\varphi \colon C \to A\) through the canonical \(f_P \colon C \to J\) as follows. The
symmetric power \(C^{(g)}\) admits a regular morphism \(\varphi^{(g)} \colon C^{(g)} \to A\)
defined by
\[
  \varphi^{(g)}\bigl(P_1 + \dots + P_g\bigr) \;=\; \varphi(P_1) + \dots + \varphi(P_g),
\]
and a birational morphism \(C^{(g)} \dashrightarrow J\) (the Abel--Jacobi map). One then
seeks a morphism \(J \to A\) making the obvious diagram commute. The construction goes via
\(C^{(g)} \dashrightarrow A\) followed by descent to \(J\); the rational map
\(C^{(g)} \dashrightarrow A\) obtained at the intermediate step is exactly the kind that
\cref{thm:rational_map_to_av_extends} promotes to a regular morphism, because \(C^{(g)}\) is
a nonsingular variety (Milne~III.3) and \(A\) is an abelian variety. Without
\cref{thm:rational_map_to_av_extends} this step has no replacement: Milne's autoduality
shortcut (\(J \cong J^\vee\)) routes through the theorem of the cube, which is not part of
the project's development; see the analysis in \cref{chap:Jacobian}.

\section{Lean encoding}
\label{sec:thm32_lean_encoding}

The Lean target file is the new
\texttt{AlgebraicJacobian/Albanese/Thm32RationalMapExtension.lean}. The single declaration
this chapter targets is the theorem
\texttt{AlgebraicGeometry.Scheme.RationalMap.extend\_to\_av}, with signature (in informal
notation):
\[
  \texttt{X : Scheme} \;\to\; \texttt{[Smooth X over \(\bar k\)]} \;\to\;
  \texttt{A : AbelianVariety \(\bar k\)} \;\to\;
  (f : \texttt{X.RationalMap A}) \;\to\;
  \exists ! \, \widetilde f \colon X \to A, \;\; \widetilde f \;\text{extends}\; f.
\]
The proof body imports \cref{thm:codim_one_extension}'s exports and
\cref{cor:regular_cohen_macaulay}'s exports from
\texttt{AlgebraicJacobian/Albanese/CodimOneExtension.lean} and
\texttt{AlgebraicJacobian/Albanese/AuslanderBuchsbaum.lean} respectively, and combines them
in the three-step pattern above (codim \(\geq 2\) from A.4.a + pure-codim \(1\) from A.4.a
\(\Rightarrow\) empty indeterminacy locus).

The Mathlib namespace target is
\texttt{AlgebraicGeometry.Scheme.RationalMap} (already housed in
\texttt{Mathlib.AlgebraicGeometry.Birational.RationalMap} at the project's pinned commit;
see \cref{chap:RiemannRoch_WeilDivisor}, where the same namespace hosts the
order-at-a-prime-divisor map). The extension is therefore a method on the existing
\texttt{Scheme.RationalMap} structure rather than a free-standing predicate, allowing
downstream consumers in \cref{chap:Albanese_AlbaneseUP} to call it as
\texttt{f.extend\_to\_av}.

No new core axiom is required; the entire chapter is a two-line combination of A.4.a and
A.4.b. The chapter is therefore inexpensive on the prover side (the LOC budget is dominated
by A.4.a and A.4.b themselves); its purpose is to record the precise statement and the
exact dependency wiring so that A.4.d can consume a single named lemma.

\section{Internal Lean helper lemmas}
\label{sec:thm32_lean_helpers}

The Lean proof of \cref{thm:rational_map_to_av_extends} factors through a small
chain of internal helper lemmas that decompose the two-input combination above
into self-contained, separately-checkable pieces. These carry no external
mathematical source; each is a packaging step internal to the formalisation.

\begin{lemma}\leanok
[A scheme smooth over an algebraically closed field is reduced]
  \label{lem:isReduced_of_smooth_over_field}
  \lean{AlgebraicGeometry.Scheme.RationalMap.isReduced_of_smooth_over_field}
  For an algebraically closed field \(\bar k\) and an object \(A\) over
  \(\Spec \bar k\) whose structure morphism \(A.\mathrm{hom}\) is smooth, the
  underlying scheme \(A.\mathrm{left}\) is reduced.
  \begin{proof}
    \uses{lem:smooth_scheme_reduced_algclosed}
    Immediate from \cref{lem:smooth_scheme_reduced_algclosed}.
  \end{proof}
\end{lemma}

\begin{lemma}\leanok
[Integrality of the abelian variety]
  \label{lem:av_isIntegral_of_smooth_geomIrred}
  \lean{AlgebraicGeometry.Scheme.RationalMap.av_isIntegral_of_smooth_geomIrred}
  \uses{lem:isReduced_of_smooth_over_field}
  For an abelian variety \(A\) over an algebraically closed field \(\bar k\)
  (encoded as \(A\) over \(\Spec \bar k\) with group-object, proper, smooth, and
  geometrically irreducible structure), the underlying scheme \(A.\mathrm{left}\)
  is integral: it is irreducible (geometric irreducibility over the singleton
  base \(\Spec \bar k\)) and reduced
  (\cref{lem:isReduced_of_smooth_over_field}).
  \begin{proof}
    Proved directly in Lean.
  \end{proof}
\end{lemma}

\begin{lemma}\leanok
[The indeterminacy locus of a rational map to an abelian variety is empty]
  \label{lem:av_indeterminacy_empty}
  \lean{AlgebraicGeometry.Scheme.RationalMap.av_indeterminacyLocus_eq_empty}
  \uses{lem:av_isIntegral_of_smooth_geomIrred, lem:milne_codim1_indeterminacy, thm:indeterminacy_codimGe2}
  For a rational map \(f \colon X \dashrightarrow A\), defined over \(\bar k\),
  from a nonsingular variety \(X\) to an abelian variety \(A\) over an
  algebraically closed field \(\bar k\), the indeterminacy locus of \(f\) is
  empty. Indeed, by the pure-codim-\(1\)-or-empty dichotomy of Milne's
  Lemma~3.3 (\cref{lem:milne_codim1_indeterminacy}, applied with
  \(A.\mathrm{left}\) integral by \cref{lem:av_isIntegral_of_smooth_geomIrred},
  since an abelian variety is in particular a group variety), the locus is
  empty or contains a point of codimension \(1\); the latter contradicts the
  codim-\(\geq 2\) conclusion of Milne's Theorem~3.1
  (\cref{thm:indeterminacy_codimGe2}).
  \begin{proof}
    Proved directly in Lean.
  \end{proof}
\end{lemma}

\begin{remark}
  \label{rmk:thm32_why_empty_not_codim1free}
  The chain deliberately runs through ``the indeterminacy locus is
  \emph{empty}'' rather than through codim-\(1\)-indeterminacy-freeness: for a
  merely complete target, codim-\(1\)-freeness does \emph{not} imply the
  existence of a regular extension (see
  \cref{rmk:codim1_no_extension_from_codim1free}). Emptiness of the locus is
  exactly the hypothesis of the extension theorem
  \cref{thm:codim_one_extension}.
\end{remark}

\section{Out of scope}
\label{sec:thm32_out_of_scope}

This chapter packages \emph{only} the theorem
\cref{thm:rational_map_to_av_extends} (Milne~3.2) and the immediate proof combining its
two inputs. The following adjacent content is explicitly out of scope here:
\begin{itemize}
  \item \textbf{The codim-\(\geq 2\) indeterminacy theorem (Milne Theorem~3.1) itself,
    and the extension theorem from an empty indeterminacy locus.} Proved in
    \cref{chap:Albanese_CodimOneExtension} as \cref{thm:indeterminacy_codimGe2}
    and \cref{thm:codim_one_extension}; this chapter only consumes them.
  \item \textbf{The pure-codim-\(1\) indeterminacy lemma (Milne Lemma~3.3) itself.}
    Proved in \cref{chap:Albanese_CodimOneExtension} as
    \cref{lem:milne_codim1_indeterminacy}; this chapter only consumes it.
  \item \textbf{The Auslander--Buchsbaum formula and its Cohen--Macaulay corollary
    (Milne~\(\equiv\) Stacks 090V / 00OD).} Proved in
    \cref{chap:Albanese_AuslanderBuchsbaum} as \cref{thm:auslander_buchsbaum} and
    \cref{cor:regular_cohen_macaulay}; the present chain no longer consumes them
    (no depth or local-cohomology input is needed once the indeterminacy locus
    is shown empty; see \cref{rmk:thm32_role_of_ab}).
  \item \textbf{Milne Theorem~3.4 / Corollary~3.7 / Theorem~3.8 / Proposition~3.9 /
    Proposition~3.10.} The remainder of Milne~\S I.3 (the special cases ``rational map
    from a product of nonsingular varieties to an AV vanishing on the axes is constant'',
    ``rational map between AVs = homomorphism + translation'', ``birationally equivalent
    AVs are isomorphic'', ``every \(\mathbb A^1 \dashrightarrow A\) is constant'',
    ``every rational map from a unirational variety to an AV is constant'') is downstream
    of \cref{thm:rational_map_to_av_extends} but is not needed by A.4 in the project's
    current routing: the genus-\(0\) base case
    \cref{rmk:base_case_fourth_route} closes
    via the proven Rigidity Lemma + \(\mathbb G_m\)-scaling shortcut without invoking any of
    these. They are recorded as off-path in
    \cref{chap:AbelianVarietyRigidity}.
  \item \textbf{The Albanese universal property of \(\Pic^0_{C/k}\) (Milne Proposition~6.1).}
    Consumed by A.4.d, proved in \cref{chap:Albanese_AlbaneseUP}; this chapter only
    provides the technical extension lemma that A.4.d invokes when promoting
    \(C^{(g)} \dashrightarrow A\) to \(C^{(g)} \to A\).
  \item \textbf{Generalisations beyond an algebraically closed base field.} Milne's
    Theorem~3.2 is stated for varieties over an algebraically closed field; the
    project's positive-genus arm runs over \(\bar k\) throughout (\texttt{chap:RigidityKbar}).
    Descent of the extension to a non-algebraically-closed base, if ever needed, would
    require additional Galois-descent machinery and is not in scope.
\end{itemize}
